\chapter{测度论与随机空间基础}

\label{chap:02}

第 1 章给出了拓扑与空间层级的语言，但“可收敛”并不自动推出“可积分”，而“可积”也不自动推出“可做随机建模”。从无穷维优化的角度看，真正进入概率、风险和动态控制之前，还需要一套能够处理事件、期望、路径和测度对偶的统一框架。本章的任务，就是把这套框架搭起来。

本章沿着四条彼此衔接的主线展开。第 \ref{sec:2-1} 节先从可测空间与测度延拓出发，解释为什么随机模型往往不是直接给出在完整 $\sigma$-代数上的测度，而是先在一个较小的可观察事件类上指定数据，再通过延拓定理补全。第 \ref{sec:2-2} 节转入 Lebesgue 积分与收敛定理，它们构成了后文“极限与期望交换”这一技术动作的合法性基础。第 \ref{sec:2-3} 节把视角推进到拓扑空间上的测度，特别是 Radon 测度与 Riesz 表示定理，这一步把抽象线性泛函首次具体化为测度对象。最后第 \ref{sec:2-4} 节在此基础上进入概率空间、滤过、鞅、Brownian motion 与 Itô 积分，为第 11 章的连续时间控制与量化建模准备最小但完整的随机分析入口。

和第 1 章一样，本章不追求百科式覆盖，而是坚持优化导向。凡是后文真正需要反复调用的概念和定理，本章都会给出标准定义、基本性质和可引用编号；凡是更深的概率论支线，例如一般半鞅理论或测度论的最广泛抽象形式，则只保留与本书主线直接相关的版本。读完这一章后，读者应当能够把“随机输入”“期望目标”“测度乘子”“连续时间噪声”这些对象放到同一张逻辑地图中，而不再把它们看成彼此无关的技术碎片。

\section{σ-代数与测度延拓}

\label{sec:2-1}

随机优化模型看上去往往以“概率分布”开始，但在严格写法里，概率分布并不是对所有子集都赋值，而是定义在某个可测结构上。更常见的情况是，我们最初只知道某些基本事件的大小，例如区间、柱集、矩形或场景树节点的权重；真正的测度则需要由这些初始数据延拓得到。本节就是要说明，从有限可观察事件到完整概率模型的这一步，如何在集合论上被严密完成。

\subsection{可测结构与生成的 $\sigma$-代数}

\begin{definition}[$\sigma$-代数]\label{def:sigma-algebra}
设 $\Omega$ 为非空集合。子集族 $\mathscr F\subseteq 2^\Omega$ 称为 $\Omega$ 上的一个\emph{$\sigma$-代数}，若满足：
\begin{enumerate}
  \item $\Omega\in \mathscr F$；
  \item 若 $A\in \mathscr F$，则 $A^c\in \mathscr F$；
  \item 若 $\{A_n\}_{n\ge 1}\subseteq \mathscr F$，则 $\bigcup_{n\ge1}A_n\in \mathscr F$。
\end{enumerate}
此时称 $(\Omega,\mathscr F)$ 为一个\emph{可测空间}，$\mathscr F$ 中的集合称为\emph{可测集}。
\end{definition}

\begin{proposition}[由集合族生成的 $\sigma$-代数]\label{prop:generated-sigma-algebra}
设 $\mathscr C\subseteq 2^\Omega$ 为任意集合族，则存在包含 $\mathscr C$ 的最小 $\sigma$-代数，记为 $\sigma(\mathscr C)$。
\end{proposition}

\begin{proof}
考虑所有包含 $\mathscr C$ 的 $\sigma$-代数之全体 $\{\mathscr G_i\}_{i\in I}$。由于幂集 $2^\Omega$ 本身是一个 $\sigma$-代数，因此这类对象非空。令
\[
\mathscr F_0:=\bigcap_{i\in I}\mathscr G_i.
\]
由 $\sigma$-代数定义逐条验证即可知，任意族 $\sigma$-代数的交仍是 $\sigma$-代数，因此 $\mathscr F_0$ 是包含 $\mathscr C$ 的 $\sigma$-代数。由构造可知，它被包含在任何其他包含 $\mathscr C$ 的 $\sigma$-代数之中，故为最小者。
\end{proof}

\begin{definition}[测度与预备测度]\label{def:measure-premeasure}
设 $(\Omega,\mathscr F)$ 为可测空间。映射 $\mu\colon \mathscr F\to [0,+\infty]$ 称为\emph{测度}，若满足
\[
\mu(\varnothing)=0
\]
且对任意两两不交的集合列 $\{A_n\}_{n\ge1}\subseteq \mathscr F$ 都有
\[
\mu\Bigl(\bigcup_{n=1}^{\infty}A_n\Bigr)=\sum_{n=1}^{\infty}\mu(A_n).
\]
若 $\mathscr A$ 只是一个代数，而 $\mu_0$ 在 $\mathscr A$ 上满足上述可列可加条件，只要并集仍留在 $\mathscr A$ 中，就称 $\mu_0$ 为 $\mathscr A$ 上的\emph{预备测度}。
\end{definition}

\begin{example}[欧氏空间中的 Borel 结构]\label{ex:borel-sigma-rn}
在 $\mathbb R^n$ 中，由所有开集生成的 $\sigma$-代数称为 Borel $\sigma$-代数，记为 $\mathcal B(\mathbb R^n)$。把这个例子放在这里，是为了说明 Boyd 书里默认出现的“约束集可测”“随机变量有分布”，实际上都隐含着定义~\ref{def:sigma-algebra} 和命题~\ref{prop:generated-sigma-algebra} 这一步已经完成。
\end{example}

\paragraph{为什么要从生成说起}

\begin{remark}
在无穷维和路径空间建模中，真正容易书写的通常不是全部可测集，而是某一类基本事件。例如在函数空间上，人们常先指定有限维投影对应的柱集；在场景树模型中，则先指定每个节点的条件概率。命题~\ref{prop:generated-sigma-algebra} 告诉我们：只要先固定好这些“原始可观察事件”，后面的测度延拓问题就有了明确的落脚点。
\end{remark}

\subsection{外测度与 Carathéodory 可测性}

\begin{definition}[外测度]\label{def:outer-measure}
设 $\Omega$ 为非空集合。映射 $\mu^\ast\colon 2^\Omega\to [0,+\infty]$ 称为 $\Omega$ 上的一个\emph{外测度}，若满足：
\begin{enumerate}
  \item $\mu^\ast(\varnothing)=0$；
  \item 若 $A\subseteq B$，则 $\mu^\ast(A)\le \mu^\ast(B)$；
  \item 对任意集合列 $\{A_n\}_{n\ge1}\subseteq 2^\Omega$，有
  \[
  \mu^\ast\Bigl(\bigcup_{n=1}^{\infty}A_n\Bigr)\le \sum_{n=1}^{\infty}\mu^\ast(A_n).
  \]
\end{enumerate}
\end{definition}

\begin{definition}[Carathéodory 可测集]\label{def:caratheodory-measurable}
设 $\mu^\ast$ 为 $\Omega$ 上的外测度。若集合 $E\subseteq \Omega$ 满足：对任意 $A\subseteq \Omega$ 都有
\[
\mu^\ast(A)=\mu^\ast(A\cap E)+\mu^\ast(A\cap E^c),
\]
则称 $E$ 为关于 $\mu^\ast$ 的\emph{Carathéodory 可测集}。
\end{definition}

\begin{theorem}[Carathéodory 可测集定理]\label{thm:caratheodory-measurable-sets}
设 $\mu^\ast$ 为 $\Omega$ 上的外测度，记
\[
\mathscr M(\mu^\ast):=\{E\subseteq \Omega:\ E\text{ 关于 }\mu^\ast\text{ 是 Carathéodory 可测的}\}.
\]
则 $\mathscr M(\mu^\ast)$ 是一个 $\sigma$-代数，并且 $\mu^\ast$ 在 $\mathscr M(\mu^\ast)$ 上的限制是一个完备测度。
\end{theorem}

\begin{proof}
先证 $\mathscr M(\mu^\ast)$ 对补运算封闭。若 $E$ 满足定义~\ref{def:caratheodory-measurable}，则对任意 $A\subseteq \Omega$，
\[
\mu^\ast(A)=\mu^\ast(A\cap E)+\mu^\ast(A\cap E^c),
\]
交换 $E$ 与 $E^c$ 的位置即可得 $E^c\in \mathscr M(\mu^\ast)$。

再证对有限交封闭。设 $E,F\in \mathscr M(\mu^\ast)$。对任意 $A\subseteq \Omega$，先按 $E$ 分割，再对与 $E$ 的交部分按 $F$ 分割，得到
\[
\mu^\ast(A)=\mu^\ast(A\cap E\cap F)+\mu^\ast(A\cap E\cap F^c)+\mu^\ast(A\cap E^c).
\]
将后两项重新组合，便可写成
\[
\mu^\ast(A)=\mu^\ast(A\cap (E\cap F))+\mu^\ast(A\cap (E\cap F)^c),
\]
故 $E\cap F\in \mathscr M(\mu^\ast)$。由补封闭也得有限并封闭。

对可数并，只需先处理两两不交的情形，再通过递增极限推广。设 $E_n\in \mathscr M(\mu^\ast)$ 两两不交，令 $E=\bigcup_{n\ge1}E_n$。由外测度的次可加性总有
\[
\mu^\ast(A)\le \mu^\ast(A\cap E)+\mu^\ast(A\cap E^c).
\]
反向不等式则由每个 $E_n$ 的可测性分步切割并取极限得到。因此 $E\in \mathscr M(\mu^\ast)$，从而 $\mathscr M(\mu^\ast)$ 为 $\sigma$-代数。

最后，若 $\{E_n\}$ 两两不交且都在 $\mathscr M(\mu^\ast)$ 中，则对 $E=\bigcup_n E_n$，利用可测性等式逐步相加可得
\[
\mu^\ast(E)=\sum_{n=1}^{\infty}\mu^\ast(E_n).
\]
故限制到 $\mathscr M(\mu^\ast)$ 后是测度。若 $\mu^\ast(N)=0$，则对任意 $A\subseteq\Omega$ 有
\[
\mu^\ast(A)\le \mu^\ast(A\cap N)+\mu^\ast(A\cap N^c)\le 0+\mu^\ast(A),
\]
所以 $N$ 的任意子集也可测，从而该测度完备。
\end{proof}

\subsection{预备测度的延拓}

\begin{theorem}[Carathéodory 测度延拓定理]\label{thm:caratheodory-extension}
设 $\mathscr A$ 为 $\Omega$ 上的代数，$\mu_0$ 为 $\mathscr A$ 上的预备测度。定义
\[
\mu^\ast(E):=\inf\Bigl\{\sum_{n=1}^{\infty}\mu_0(A_n):\ E\subseteq \bigcup_{n=1}^{\infty}A_n,\ A_n\in \mathscr A\Bigr\},\qquad E\subseteq \Omega.
\]
则：
\begin{enumerate}
  \item $\mu^\ast$ 是 $\Omega$ 上的外测度；
  \item $\mathscr A\subseteq \mathscr M(\mu^\ast)$；
  \item $\mu:=\mu^\ast|_{\sigma(\mathscr A)}$ 是 $\sigma(\mathscr A)$ 上的测度，且 $\mu|_{\mathscr A}=\mu_0$；
  \item 若 $\mu_0$ 是 $\sigma$-有限的，则该延拓在 $\sigma(\mathscr A)$ 上唯一。
\end{enumerate}
\end{theorem}

\begin{proof}
由定义可直接验证 $\mu^\ast(\varnothing)=0$ 与单调性。次可加性来自于对每个 $E_n$ 选取近似最优覆盖后，把这些覆盖拼接起来。

下证 $\mathscr A\subseteq \mathscr M(\mu^\ast)$。任取 $E\in \mathscr A$ 和任意 $B\subseteq \Omega$。对覆盖 $B\subseteq \bigcup_n A_n$，由代数对交、补封闭可知 $A_n\cap E$ 与 $A_n\cap E^c$ 仍在 $\mathscr A$ 中，于是
\[
\mu_0(A_n)=\mu_0(A_n\cap E)+\mu_0(A_n\cap E^c).
\]
对所有覆盖取下确界即可得到
\[
\mu^\ast(B)\ge \mu^\ast(B\cap E)+\mu^\ast(B\cap E^c).
\]
反向不等式由外测度的次可加性总是成立，因此 $E$ 可测。

由定理~\ref{thm:caratheodory-measurable-sets}，$\mu^\ast$ 在 $\mathscr M(\mu^\ast)$ 上是测度；既然 $\sigma(\mathscr A)\subseteq \mathscr M(\mu^\ast)$，便得到了所需延拓。对于 $E\in \mathscr A$，以自身作为覆盖可得 $\mu^\ast(E)\le \mu_0(E)$；反向不等式由上一步切割及预备测度的可加性可得，所以 $\mu^\ast(E)=\mu_0(E)$。

若 $\mu_0$ 是 $\sigma$-有限的，则可取 $\Omega=\bigcup_n E_n$，其中 $E_n\in \mathscr A$ 且 $\mu_0(E_n)<\infty$。若 $\nu$ 是另一延拓，则 $\mu$ 与 $\nu$ 在生成代数 $\mathscr A$ 上一致，并在每个有限测度块 $E_n$ 上都有限。由单调类论证可知二者在 $\sigma(\mathscr A)$ 上一致，从而唯一。
\end{proof}

\begin{example}[由分布函数生成测度]\label{ex:distribution-function-measure}
设 $F\colon \mathbb R\to \mathbb R$ 为右连续、单调不减函数。若在半开区间代数上规定
\[
\mu_0((a,b])=F(b)-F(a),
\]
则 $\mu_0$ 是一个预备测度。由定理~\ref{thm:caratheodory-extension}，它唯一延拓为 $\mathcal B(\mathbb R)$ 上的 Borel 测度。把这个例子放在这里，是为了说明随机变量的分布函数并不是“额外数据”，而正是一个通过延拓定理生成测度的起点。
\end{example}

\begin{example}[抽象例子：路径空间上的柱集]\label{ex:cylinder-extension}
在路径空间 $\mathbb R^{[0,T]}$ 或 $C([0,T])$ 上，最容易指定概率的对象往往是有限维柱集，即只依赖有限个时刻取值的事件。先在柱集代数上指定一致的有限维分布，再通过延拓定理得到完整路径测度，是构造 Brownian motion 和更一般随机过程的标准路线。把这个例子放在这里，是为了给第 \ref{sec:2-4} 节的随机过程语言留出严格基础。
\end{example}

\begin{example}[工程触点：场景树到概率空间]\label{ex:scenario-tree-measure}
在供应链或随机规划中，人们常先给出有限层场景树，每个节点带有转移概率。若要在极限模型中讨论期望目标、风险约束或动态一致性，就需要把这些离散层级数据嵌入一个真正的概率空间。把这个例子放在这里，是为了说明测度延拓并非纯粹抽象技巧，而是从有限场景模型走向路径级随机优化的第一步。
\end{example}

\subsection*{有限维对照}

有限维优化里，随机输入通常直接写成“随机向量 $X\in \mathbb R^n$ 有某个分布”，因此底层的可测结构和测度延拓常被隐藏起来。本节的作用，是把这些被默认完成的步骤显式写出：从定义~\ref{def:sigma-algebra} 到定义~\ref{def:outer-measure}，再到定理~\ref{thm:caratheodory-extension}。

\subsection*{习题（待补）}

\begin{exercise}[分布函数延拓占位]\label{exr:distribution-extension-placeholder}
补一题：从例~\ref{ex:distribution-function-measure} 出发，证明给定右连续分布函数 $F$ 后，区间赋值 $(a,b]\mapsto F(b)-F(a)$ 的确构成一个预备测度。
\end{exercise}

\begin{exercise}[柱集桥接题占位]\label{exr:cylinder-extension-placeholder}
补一题：在有限维情形把柱集代数退化为矩形代数，说明为什么定理~\ref{thm:caratheodory-extension} 正是“联合分布存在”的集合论版本。
\end{exercise}


\section{积分论与收敛定理}

\label{sec:2-2}

测度只给出了“集合有多大”，真正进入优化与概率时，我们关心的是函数在这些集合上的平均、总量与极限行为。Lebesgue 积分之所以在本书中不可替代，并不只是因为它比 Riemann 积分更一般，而是因为它允许我们在极限过程里保持对可积性的控制。本节的核心，就是把“积分是什么”和“什么时候可以把极限推进积分号内”两件事组织成一条可反复调用的技术链。

\subsection{Lebesgue 积分的构造}

\begin{definition}[简单函数与非负积分]\label{def:simple-function-integral}
设 $(\Omega,\mathscr F,\mu)$ 为测度空间。若函数 $\varphi\colon \Omega\to [0,+\infty)$ 可写成
\[
\varphi=\sum_{k=1}^{m} a_k \mathbf 1_{A_k},
\]
其中 $a_k\ge 0$，$A_k\in \mathscr F$ 两两不交，则称 $\varphi$ 为一个\emph{非负简单函数}，并定义
\[
\int_\Omega \varphi\,d\mu:=\sum_{k=1}^{m} a_k\,\mu(A_k).
\]
\end{definition}

\begin{definition}[Lebesgue 积分]\label{def:lebesgue-integral}
设 $f\colon \Omega\to [0,+\infty]$ 为非负可测函数。定义
\[
\int_\Omega f\,d\mu
:=
\sup\Bigl\{\int_\Omega \varphi\,d\mu:\ 0\le \varphi\le f,\ \varphi\text{ 为简单函数}\Bigr\}.
\]
对一般实值可测函数 $f$，写
\[
f^+=\max\{f,0\},\qquad f^-=\max\{-f,0\}.
\]
若
\[
\int_\Omega f^+\,d\mu<\infty,\qquad \int_\Omega f^-\,d\mu<\infty,
\]
则称 $f$ \emph{可积}，并定义
\[
\int_\Omega f\,d\mu:=\int_\Omega f^+\,d\mu-\int_\Omega f^-\,d\mu.
\]
\end{definition}

\paragraph{为什么强调 $L^1$}

\begin{remark}
定义~\ref{def:lebesgue-integral} 给出的可积函数类正是 $L^1(\mu)$ 的底层对象。后文讨论期望目标、风险泛函、弱收敛下的积分连续性时，真正起作用的往往不是函数的逐点值，而是它在 $L^1$ 中的等价类与积分控制。
\end{remark}

\begin{proposition}[积分的基本性质]\label{prop:integral-basic-properties}
若 $f,g$ 可积，$a,b\in \mathbb R$，则：
\begin{enumerate}
  \item $af+bg$ 可积，且
  \[
  \int (af+bg)\,d\mu=a\int f\,d\mu+b\int g\,d\mu;
  \]
  \item 若 $f\le g$ a.e.，则
  \[
  \int f\,d\mu\le \int g\,d\mu;
  \]
  \item
  \[
  \left|\int f\,d\mu\right|\le \int |f|\,d\mu.
  \]
\end{enumerate}
\end{proposition}

\begin{proof}
对简单函数结论显然，由上确界定义与单调逼近推广到非负可测函数；再对一般可积函数分解为正负部分即可。第三条由
\[
-|f|\le f\le |f|
\]
和第二条直接推出。
\end{proof}

\subsection{单调收敛、Fatou 引理与控制收敛}

\begin{theorem}[单调收敛定理]\label{thm:monotone-convergence}
设 $\{f_n\}_{n\ge1}$ 为非负可测函数列，且
\[
0\le f_1\le f_2\le \cdots,\qquad f_n\uparrow f \text{ a.e.}
\]
则
\[
\int_\Omega f_n\,d\mu \uparrow \int_\Omega f\,d\mu.
\]
\end{theorem}

\begin{proof}
由单调性先得
\[
\int f_n\,d\mu\le \int f\,d\mu,
\]
故左侧极限存在且不超过右侧。设
\[
L:=\sup_n \int f_n\,d\mu.
\]
任取简单函数 $\varphi$ 满足 $0\le \varphi\le f$。对任意 $\varepsilon\in(0,1)$，考虑集合
\[
E_n:=\{\,\omega:\ f_n(\omega)\ge (1-\varepsilon)\varphi(\omega)\,\}.
\]
由于 $f_n\uparrow f\ge \varphi$，故 $\mathbf 1_{E_n}\uparrow 1$。于是由简单函数积分与测度从下连续性可得
\[
\int (1-\varepsilon)\varphi\,d\mu
=
\lim_{n\to\infty}\int_{E_n}(1-\varepsilon)\varphi\,d\mu
\le
\limsup_{n\to\infty}\int f_n\,d\mu
=L.
\]
令 $\varepsilon\downarrow 0$，再对所有 $\varphi\le f$ 取上确界，得到
\[
\int f\,d\mu\le L.
\]
结合前面的反向不等式，结论成立。
\end{proof}

\begin{lemma}[Fatou 引理]\label{lem:fatou}
设 $\{f_n\}_{n\ge1}$ 为非负可测函数列，则
\[
\int_\Omega \liminf_{n\to\infty} f_n\,d\mu
\le
\liminf_{n\to\infty}\int_\Omega f_n\,d\mu.
\]
\end{lemma}

\begin{proof}
设
\[
g_n:=\inf_{k\ge n} f_k.
\]
则 $\{g_n\}$ 仍为非负可测函数列，且 $g_n\uparrow \liminf_{n\to\infty}f_n$。由定理~\ref{thm:monotone-convergence}，
\[
\int \liminf_{n\to\infty}f_n\,d\mu
=
\lim_{n\to\infty}\int g_n\,d\mu.
\]
又对每个固定 $n$ 与任意 $k\ge n$，都有 $g_n\le f_k$，故
\[
\int g_n\,d\mu\le \inf_{k\ge n}\int f_k\,d\mu.
\]
令 $n\to\infty$ 即得结论。
\end{proof}

\begin{theorem}[控制收敛定理]\label{thm:dominated-convergence}
设 $\{f_n\}_{n\ge1}$ 为可测函数列，且
\[
f_n\to f \text{ a.e.},\qquad |f_n|\le g \text{ a.e.},
\]
其中 $g$ 可积。则 $f$ 可积，且
\[
\int_\Omega f_n\,d\mu\to \int_\Omega f\,d\mu.
\]
\end{theorem}

\begin{proof}
由 $|f_n|\le g$ 和逐点收敛可得 $|f|\le g$ a.e.，故 $f$ 可积。对非负函数列 $g+f_n$ 应用 Fatou 引理，得
\[
\int (g+f)\,d\mu
\le
\liminf_{n\to\infty}\int (g+f_n)\,d\mu
=
\int g\,d\mu+\liminf_{n\to\infty}\int f_n\,d\mu.
\]
从而
\[
\int f\,d\mu\le \liminf_{n\to\infty}\int f_n\,d\mu.
\]
同理，对非负函数列 $g-f_n$ 应用 Fatou 引理，得到
\[
\int (g-f)\,d\mu
\le
\liminf_{n\to\infty}\int (g-f_n)\,d\mu
=
\int g\,d\mu-\limsup_{n\to\infty}\int f_n\,d\mu,
\]
即
\[
\limsup_{n\to\infty}\int f_n\,d\mu\le \int f\,d\mu.
\]
结合二者便得所求极限。
\end{proof}

\begin{proposition}[积分的绝对连续性]\label{prop:absolute-continuity-integral}
若 $f\in L^1(\mu)$，则对任意 $\varepsilon>0$，存在 $\delta>0$，使得只要 $A\in \mathscr F$ 且 $\mu(A)<\delta$，就有
\[
\int_A |f|\,d\mu<\varepsilon.
\]
\end{proposition}

\begin{proof}
取 $M>0$ 使
\[
\int_{\{|f|>M\}} |f|\,d\mu<\frac{\varepsilon}{2}.
\]
再令 $\delta=\varepsilon/(2M)$。若 $\mu(A)<\delta$，则
\[
\int_A |f|\,d\mu
=
\int_{A\cap\{|f|\le M\}}|f|\,d\mu+\int_{A\cap\{|f|>M\}}|f|\,d\mu
\le
M\mu(A)+\frac{\varepsilon}{2}
<
\varepsilon.
\]
\end{proof}

\begin{example}[抽象例子：指标函数逼近]\label{ex:indicator-approximation}
设 $A_n\uparrow A$，则 $\mathbf 1_{A_n}\uparrow \mathbf 1_A$。由定理~\ref{thm:monotone-convergence} 得
\[
\mu(A_n)=\int \mathbf 1_{A_n}\,d\mu \uparrow \int \mathbf 1_A\,d\mu=\mu(A).
\]
把这个例子放在这里，是为了说明测度从下连续性其实是积分理论中最基本的单调收敛现象。
\end{example}

\begin{example}[有限维坍缩：随机向量的期望]\label{ex:finite-dimensional-expectation}
若 $X$ 是 $\mathbb R^n$ 值随机向量，$h\colon \mathbb R^n\to \mathbb R$ 可积，则
\[
\mathbb E[h(X)]=\int_{\mathbb R^n} h(x)\,\mu_X(dx).
\]
把这个例子放在这里，是为了提醒读者 Boyd 书里常见的“期望目标函数”并不是一种独立算子，而只是定义~\ref{def:lebesgue-integral} 在概率测度上的写法。
\end{example}

\begin{example}[应用触点：Monte Carlo 与风险泛函]\label{ex:monte-carlo-risk}
设损失函数 $L_n$ 由模拟近似得到，且 $L_n\to L$ a.s.，同时存在可积上界 $G$ 使 $|L_n|\le G$。由定理~\ref{thm:dominated-convergence}，
\[
\mathbb E[L_n]\to \mathbb E[L].
\]
把这个例子放在这里，是为了说明后文量化风险优化和分布鲁棒目标中，控制收敛定理是“极限可推进期望号内”的最基本合法性来源。
\end{example}

\subsection*{有限维对照}

有限维教材里，人们经常把“对随机变量取期望”“把极限与期望交换”写成直观计算；而在无穷维和函数空间中，这些动作是否成立完全依赖定理~\ref{thm:monotone-convergence}、引理~\ref{lem:fatou} 和定理~\ref{thm:dominated-convergence} 所提供的条件链。

\subsection*{习题（待补）}

\begin{exercise}[单调收敛占位]\label{exr:monotone-convergence-placeholder}
补一题：用定理~\ref{thm:monotone-convergence} 证明概率测度对递增事件列的从下连续性。
\end{exercise}

\begin{exercise}[控制收敛桥接题占位]\label{exr:dominated-convergence-placeholder}
补一题：把例~\ref{ex:monte-carlo-risk} 落回欧氏空间中的随机向量积分，说明控制上界在数值优化中为何常以矩条件或截断条件的形式出现。
\end{exercise}


\section{拓扑空间上的测度}

\label{sec:2-3}

前两节给出的测度与积分理论，还停留在纯集合论层面；但在优化、概率和对偶理论中，我们常常需要让测度与拓扑同时出现。一个典型场景是：在某个状态空间上先有连续函数，再考虑“对所有连续测试函数积分”所定义的线性泛函。若要把这种泛函解释成真正的测度，就必须把拓扑结构纳入讨论。对本书而言，最重要的环境是局部紧 Hausdorff 空间，以及第 \ref{sec:1-3} 节已经建立的 Polish 空间框架；前者适合陈述经典 Riesz 表示定理，后者则保证了许多在随机分析里需要的可分性与正则性。

\subsection{Borel 测度与 Radon 测度}

设 $X$ 为拓扑空间，记其 Borel $\sigma$-代数为 $\mathcal B(X)$。定义在 $\mathcal B(X)$ 上的测度称为 $X$ 上的 \emph{Borel 测度}。仅有 Borel 性通常还不够，因为它并不能保证测度与拓扑上的紧集、开集良好配合；真正适合分析与优化的，是正则性更强的 Radon 测度。

\begin{definition}[Radon 测度]\label{def:radon-measure}
设 $X$ 为局部紧 Hausdorff 空间。定义在 $\mathcal B(X)$ 上的 Borel 测度 $\mu$ 称为一个\emph{Radon 测度}，若满足：
\begin{enumerate}
  \item 对每个紧集 $K\subseteq X$，都有 $\mu(K)<\infty$；
  \item 对每个开集 $U\subseteq X$，有
  \[
  \mu(U)=\sup\{\mu(K):\ K\subseteq U,\ K\text{ 为紧集}\};
  \]
  \item 对每个 Borel 集 $A\subseteq X$，有
  \[
  \mu(A)=\inf\{\mu(U):\ A\subseteq U,\ U\text{ 为开集}\}.
  \]
\end{enumerate}
\end{definition}

\paragraph{为何强调第 \ref{sec:1-3} 节的 Polish 环境}

\begin{remark}
若 $X$ 是定义~\ref{def:polish-space} 中的 Polish 空间，则其 Borel 结构具有很强的稳定性；在许多自然情形下，有限 Borel 测度自动具有良好的正则性。换言之，第 \ref{sec:1-3} 节并不是与测度论平行的一条支线，而是在为本节“测度与拓扑相容”提供常用环境。
\end{remark}

\begin{proposition}[Radon 测度的正则性]\label{prop:radon-regularity}
设 $X$ 为局部紧 Hausdorff 空间，$\mu$ 为 $X$ 上的 Radon 测度。则对每个 Borel 集 $A\subseteq X$ 都有
\[
\mu(A)=\sup\{\mu(K):\ K\subseteq A,\ K\text{ 为紧集}\}
=\inf\{\mu(U):\ A\subseteq U,\ U\text{ 为开集}\}.
\]
特别地，对任意 $\varepsilon>0$，存在紧集 $K$ 和开集 $U$ 使
\[
K\subseteq A\subseteq U,\qquad \mu(U\setminus K)<\varepsilon.
\]
\end{proposition}

\begin{proof}
外正则性已包含在定义~\ref{def:radon-measure} 中。下面说明内正则性可由开集内正则性推出。设
\[
\mathscr M:=\bigl\{A\in \mathcal B(X):\ \mu(A)=\sup\{\mu(K):\ K\subseteq A,\ K\text{ 为紧集}\}\bigr\}.
\]
由定义可知所有开集都属于 $\mathscr M$。再利用测度对递增并与递减交的连续性，可以验证 $\mathscr M$ 对差集的适当操作和单调极限稳定，因此形成一个包含开集的单调类。由于开集生成 $\mathcal B(X)$，由单调类定理可知 $\mathscr M=\mathcal B(X)$，于是任意 Borel 集都内正则。

最后，给定 $\varepsilon>0$，先由外正则性取开集 $U\supseteq A$ 使
\[
\mu(U)\le \mu(A)+\frac{\varepsilon}{2}.
\]
再由开集的内正则性取紧集 $K\subseteq A$ 使
\[
\mu(K)\ge \mu(A)-\frac{\varepsilon}{2}.
\]
于是
\[
\mu(U\setminus K)=\mu(U)-\mu(K)\le \varepsilon,
\]
结论成立。
\end{proof}

\subsection{连续函数上的正线性泛函}

设 $X$ 为局部紧 Hausdorff 空间，记 $C_c(X)$ 为所有紧支撑连续实函数组成的线性空间，记 $C_0(X)$ 为在无穷远处消失的连续函数空间。对优化与对偶理论而言，正线性泛函最值得关注：它们在函数空间对偶中扮演“非负乘子”的角色。

\begin{definition}[正线性泛函]
设 $L$ 是定义在 $C_c(X)$ 或 $C_0(X)$ 上的线性泛函。若对任意满足 $f\ge 0$ 的函数 $f$ 都有
\[
L(f)\ge 0,
\]
则称 $L$ 为\emph{正线性泛函}。
\end{definition}

\begin{lemma}[由正泛函诱导的集合函数]\label{lem:positive-functional-contents}
设 $L$ 为 $C_c(X)$ 上的正线性泛函。对开集 $U\subseteq X$ 定义
\[
\nu(U):=\sup\{L(f):\ f\in C_c(X),\ 0\le f\le 1,\ \operatorname{supp}(f)\subseteq U\}.
\]
则 $\nu$ 在开集族上单调，且对递增开集列满足
\[
\nu\Bigl(\bigcup_{n=1}^{\infty}U_n\Bigr)=\lim_{n\to\infty}\nu(U_n).
\]
\end{lemma}

\begin{proof}
单调性由定义直接得出。若 $U_n\uparrow U$，则显然 $\nu(U_n)\le \nu(U)$，故
\[
\limsup_{n\to\infty}\nu(U_n)\le \nu(U).
\]
反过来，任取 $f\in C_c(X)$，满足 $0\le f\le1$ 且 $\operatorname{supp}(f)\subseteq U$。由于 $\operatorname{supp}(f)$ 为紧集，而 $\{U_n\}$ 递增覆盖 $U$，存在某个 $N$ 使 $\operatorname{supp}(f)\subseteq U_N$。故
\[
L(f)\le \nu(U_N)\le \sup_n \nu(U_n).
\]
对所有这类 $f$ 取上确界，得到
\[
\nu(U)\le \sup_n \nu(U_n)=\lim_{n\to\infty}\nu(U_n).
\]
二者合并即得结论。
\end{proof}

\subsection{Riesz 表示定理}

\begin{theorem}[Riesz 表示定理]\label{thm:riesz-representation}
设 $X$ 为局部紧 Hausdorff 空间，$L$ 为 $C_c(X)$ 上的正线性泛函。则存在唯一的 Radon 测度 $\mu$，使得对任意 $f\in C_c(X)$ 都有
\[
L(f)=\int_X f\,d\mu.
\]
若进一步把 $L$ 看作 $C_0(X)$ 上的有界正线性泛函，则同一结论成立，且
\[
\|L\|=\mu(X)
\]
在 $\mu(X)<\infty$ 时有效。
\end{theorem}

\begin{proof}
证明分三步进行。

第一步，由引理~\ref{lem:positive-functional-contents} 在开集上构造集合函数 $\nu$。再对任意集合 $A\subseteq X$ 定义
\[
\mu^\ast(A):=\inf\{\nu(U):\ A\subseteq U,\ U\text{ 为开集}\}.
\]
由定义可知 $\mu^\ast(\varnothing)=0$ 且单调。利用开集的可数覆盖与引理~\ref{lem:positive-functional-contents} 的递增连续性，可以验证 $\mu^\ast$ 是一个外测度。

第二步，对 $\mu^\ast$ 应用定理~\ref{thm:caratheodory-measurable-sets}。由于开集在构造中扮演基础角色，而连续函数可用局部紧 Hausdorff 空间上的 Urysohn 型截断函数逼近，能够验证所有开集都是 Carathéodory 可测的，从而 $\mathcal B(X)\subseteq \mathscr M(\mu^\ast)$。于是 $\mu:=\mu^\ast|_{\mathcal B(X)}$ 成为一个 Borel 测度。再由定义可知 $\mu$ 对开集外正则；而由 $\nu$ 的定义与紧支撑函数逼近，可以推出对开集的内正则和对紧集的有限性，因此 $\mu$ 是 Radon 测度。

第三步，证明积分表示。对满足 $0\le f\le1$ 的紧支撑连续函数，先用层析逼近和简单函数近似比较 $L(f)$ 与 $\int f\,d\mu$，得到二者相等；再由线性性推广到一般实值 $f\in C_c(X)$。唯一性来自这样的事实：若两个 Radon 测度对所有 $f\in C_c(X)$ 的积分都相等，则它们对开集的取值相同，从而由命题~\ref{prop:radon-regularity} 对所有 Borel 集都相同。
\end{proof}

\paragraph{“泛函 = 测度”的第一座桥}

\begin{remark}
定理~\ref{thm:riesz-representation} 说明，连续函数空间上的正泛函并不是神秘的抽象对偶元素，而可以唯一解释为某个 Radon 测度的积分算子。这一步会在后文至少以三种形式回归：其一是在对偶理论里把正约束乘子解释成测度；其二是在最优控制里把轨道分布、占据测度与线性泛函联系起来；其三是在概率论中把期望写成对概率测度的积分。
\end{remark}

\begin{example}[抽象例子：Dirac 测度与点值泛函]\label{ex:dirac-functional}
对任意固定 $x_0\in X$，定义
\[
L_{x_0}(f):=f(x_0),\qquad f\in C_c(X).
\]
这是一个正线性泛函。由定理~\ref{thm:riesz-representation}，它对应的 Radon 测度正是 Dirac 测度 $\delta_{x_0}$。把这个例子放在这里，是为了说明最简单的“泛函 = 测度”并不是平均，而是点值本身；这也是后文极点、分布和采样算子的基本原型。
\end{example}

\begin{example}[有限维坍缩：欧氏空间中的积分泛函]\label{ex:euclidean-radon}
取 $X=\mathbb R^n$。若 $g\in L^1_{\mathrm{loc}}(\mathbb R^n)$ 且 $g\ge0$，则
\[
L(f):=\int_{\mathbb R^n} f(x)g(x)\,dx,\qquad f\in C_c(\mathbb R^n),
\]
定义了一个正线性泛函。这里对应的 Radon 测度就是 $\mu(dx)=g(x)\,dx$。把这个例子放在这里，是为了说明有限维教材中最熟悉的“加权积分”其实已经是 Riesz 表示定理的特殊情形。
\end{example}

\begin{example}[应用触点：测度型乘子与占据测度]\label{ex:occupation-measure}
在连续时间控制与分布鲁棒优化中，约束往往不是有限个标量不等式，而是对一族测试函数同时成立的积分关系。此时拉格朗日乘子自然不再是向量，而是某种 Radon 测度。把这个例子放在这里，是为了说明定理~\ref{thm:riesz-representation} 并非只服务于纯分析，它直接决定了后文 KKT 条件里乘子的对象类型。
\end{example}

\subsection*{有限维对照}

在有限维空间里，人们常把线性泛函写成向量内积，把概率或权重写成密度函数，于是“泛函等于积分”这件事显得理所当然。本节把这种直觉放回一般拓扑空间中，并用定义~\ref{def:radon-measure}、命题~\ref{prop:radon-regularity} 和定理~\ref{thm:riesz-representation} 说明：只有当测度具有足够的正则性时，这种解释才真正稳固。

\subsection*{习题（待补）}

\begin{exercise}[Dirac 测度占位]\label{exr:dirac-riesz-placeholder}
补一题：证明例~\ref{ex:dirac-functional} 中的点值泛函对应的确是 Dirac 测度，并据此验证定理~\ref{thm:riesz-representation} 的唯一性在该情形下如何显现。
\end{exercise}

\begin{exercise}[有限维桥接题占位]\label{exr:euclidean-riesz-placeholder}
补一题：在 $X=\mathbb R^n$ 上，从例~\ref{ex:euclidean-radon} 出发验证命题~\ref{prop:radon-regularity}，说明 Lebesgue 测度为何是一个 Radon 测度。
\end{exercise}


\section{概率空间与随机过程微积分初步}

\label{sec:2-4}

前面三节已经准备好了事件、积分与拓扑测度的语言，但连续时间优化还需要进一步回答三个问题：随机对象怎样随时间演化，信息怎样逐步显露，噪声如何进入动力系统。为此，本节从概率空间与随机变量出发，经过条件期望、滤过、鞅与停时，最终进入 Brownian motion、Itô 积分与 Itô 公式。目标不是写成一部完整随机分析专著，而是建立后续最优控制、量化交易与路径空间优化真正会调用的最小语法。

\subsection{概率空间、随机变量与条件期望}

\begin{definition}[概率空间]\label{def:probability-space}
三元组 $(\Omega,\mathscr F,\mathbb P)$ 称为一个\emph{概率空间}，若 $(\Omega,\mathscr F)$ 是可测空间，且 $\mathbb P$ 是满足
\[
\mathbb P(\Omega)=1
\]
的测度。定义在 $(\Omega,\mathscr F)$ 上、取值于某个可测空间 $(E,\mathscr E)$ 的可测映射 $X\colon \Omega\to E$ 称为一个\emph{随机变量}。
\end{definition}

\paragraph{期望只是积分}

\begin{remark}
一旦概率空间被固定，随机变量的期望只是第 \ref{sec:2-2} 节 Lebesgue 积分的特殊记号：
\[
\mathbb E[X]=\int_\Omega X\,d\mathbb P.
\]
因此所有与极限交换有关的合法性，本质上都回到定理~\ref{thm:monotone-convergence}、引理~\ref{lem:fatou} 与定理~\ref{thm:dominated-convergence}。
\end{remark}

\begin{definition}[条件期望]\label{def:conditional-expectation}
设 $(\Omega,\mathscr F,\mathbb P)$ 为概率空间，$\mathscr G\subseteq \mathscr F$ 为子 $\sigma$-代数，$X\in L^1(\mathbb P)$。若随机变量 $Y$ 满足：
\begin{enumerate}
  \item $Y$ 是 $\mathscr G$-可测的；
  \item 对任意 $A\in \mathscr G$，都有
  \[
  \int_A Y\,d\mathbb P=\int_A X\,d\mathbb P,
  \]
\end{enumerate}
则称 $Y$ 为 $X$ 关于 $\mathscr G$ 的\emph{条件期望}，记为
\[
Y=\mathbb E[X\mid \mathscr G].
\]
\end{definition}

\begin{proposition}[条件期望的基本性质]\label{prop:conditional-expectation-basic}
设 $X,Y\in L^1(\mathbb P)$，$\mathscr G\subseteq \mathscr F$ 为子 $\sigma$-代数，则：
\begin{enumerate}
  \item $\mathbb E[X\mid \mathscr G]$ 在几乎处处意义下唯一；
  \item 条件期望关于 $X$ 线性；
  \item 若 $X$ 本身 $\mathscr G$-可测，则
  \[
  \mathbb E[X\mid \mathscr G]=X \qquad \text{a.s.};
  \]
  \item 若 $X\ge 0$ a.s.，则
  \[
  \mathbb E[X\mid \mathscr G]\ge 0 \qquad \text{a.s.};
  \]
  \item 塔性质成立：若 $\mathscr H\subseteq \mathscr G\subseteq \mathscr F$，则
  \[
  \mathbb E[\mathbb E[X\mid \mathscr G]\mid \mathscr H]=\mathbb E[X\mid \mathscr H]
  \qquad \text{a.s.}
  \]
\end{enumerate}
\end{proposition}

\begin{proof}
唯一性来自这样一个事实：若 $Y_1,Y_2$ 都满足定义~\ref{def:conditional-expectation}，则对任意 $A\in \mathscr G$ 有
\[
\int_A (Y_1-Y_2)\,d\mathbb P=0.
\]
取 $A=\{Y_1>Y_2\}\in \mathscr G$ 与 $A=\{Y_2>Y_1\}$ 即得 $Y_1=Y_2$ a.s.。线性与正性由积分恒等式直接验证。若 $X$ 是 $\mathscr G$-可测，则它自己就满足定义中的两条。塔性质只需对任意 $A\in \mathscr H$ 比较积分：
\[
\int_A \mathbb E[\mathbb E[X\mid \mathscr G]\mid \mathscr H]\,d\mathbb P
=\int_A \mathbb E[X\mid \mathscr G]\,d\mathbb P
=\int_A X\,d\mathbb P.
\]
故该随机变量满足 $\mathbb E[X\mid \mathscr H]$ 的刻画。
\end{proof}

\subsection{独立性、滤过、适应过程与鞅}

\begin{definition}[独立性、滤过与适应性]
设 $(\Omega,\mathscr F,\mathbb P)$ 为概率空间。
\begin{enumerate}
  \item 两个子 $\sigma$-代数 $\mathscr G_1,\mathscr G_2\subseteq \mathscr F$ 称为\emph{独立}的，若对任意 $A\in \mathscr G_1$、$B\in \mathscr G_2$，有
  \[
  \mathbb P(A\cap B)=\mathbb P(A)\mathbb P(B).
  \]
  \item 一族子 $\sigma$-代数 $\{\mathscr F_t\}_{t\ge0}$ 若满足 $s\le t$ 时
  \[
  \mathscr F_s\subseteq \mathscr F_t,
  \]
  则称其为一个\emph{滤过}。
  \item 过程 $X=\{X_t\}_{t\ge0}$ 若对每个 $t$，$X_t$ 都是 $\mathscr F_t$-可测的，则称 $X$ 对该滤过\emph{适应}。
\end{enumerate}
\end{definition}

\begin{definition}[鞅]\label{def:martingale}
设 $(\Omega,\mathscr F,\mathbb P)$ 配有滤过 $\{\mathscr F_t\}_{t\ge0}$。若过程 $M=\{M_t\}_{t\ge0}$ 满足：
\begin{enumerate}
  \item $M$ 适应于 $\{\mathscr F_t\}$；
  \item 对每个 $t$，有 $M_t\in L^1(\mathbb P)$；
  \item 对任意 $0\le s\le t$，有
  \[
  \mathbb E[M_t\mid \mathscr F_s]=M_s \qquad \text{a.s.},
  \]
\end{enumerate}
则称 $M$ 为一个\emph{鞅}。若上式中的等号分别改成“$\ge$”或“$\le$”，则得到次鞅与上鞅。
\end{definition}

\begin{definition}[停时]
设给定滤过 $\{\mathscr F_t\}_{t\ge0}$。随机变量 $\tau\colon \Omega\to [0,\infty]$ 称为一个\emph{停时}，若对每个 $t\ge0$，
\[
\{\tau\le t\}\in \mathscr F_t.
\]
\end{definition}

\begin{proposition}[停时与鞅的基本前瞻]\label{prop:optional-sampling-preview}
设 $M$ 为离散时间鞅，$\tau$ 为有界停时，则停时过程 $\{M_{n\wedge \tau}\}$ 仍为鞅，因而
\[
\mathbb E[M_\tau]=\mathbb E[M_0].
\]
\end{proposition}

\begin{proof}
对每个 $n$，变量 $M_{n\wedge \tau}$ 对 $\mathscr F_n$ 可测且可积。再由分解
\[
M_{(n+1)\wedge \tau}=M_n\mathbf 1_{\{\tau\le n\}}+M_{n+1}\mathbf 1_{\{\tau>n\}}
\]
并对 $\mathscr F_n$ 取条件期望，利用 $\{\tau>n\}\in \mathscr F_n$ 和定义~\ref{def:martingale} 的鞅性质即可得到
\[
\mathbb E[M_{(n+1)\wedge \tau}\mid \mathscr F_n]=M_{n\wedge \tau}.
\]
从而停时过程仍为鞅。再取期望便得结论。
\end{proof}

\subsection{Brownian motion 与 Itô 积分}

\begin{definition}[Brownian motion]\label{def:brownian-motion}
设 $(\Omega,\mathscr F,\mathbb P)$ 为概率空间，过程 $W=\{W_t\}_{t\ge0}$ 称为一个\emph{Brownian motion}，若满足：
\begin{enumerate}
  \item $W_0=0$ a.s.；
  \item $W$ 具有独立增量：对任意 $0\le t_0<t_1<\cdots<t_n$，随机变量
  \[
  W_{t_1}-W_{t_0},\dots,W_{t_n}-W_{t_{n-1}}
  \]
  相互独立；
  \item $W_t-W_s\sim N(0,t-s)$ 对任意 $0\le s<t$ 成立；
  \item 样本路径几乎处处连续。
\end{enumerate}
\end{definition}

\begin{example}[有限维高斯坍缩]\label{ex:gaussian-vector}
对任意 $0<t_1<\cdots<t_n$，向量
\[
(W_{t_1},\dots,W_{t_n})
\]
是中心高斯向量，其协方差矩阵为 $(\min\{t_i,t_j\})_{i,j}$。把这个例子放在这里，是为了说明 Brownian motion 可以看成一族彼此相容的有限维高斯向量在时间上的连续拼接。
\end{example}

\begin{definition}[简单适应过程上的 Itô 积分]\label{def:ito-integral}
设 $W$ 是适应于滤过 $\{\mathscr F_t\}$ 的 Brownian motion，$T>0$。若过程 $H$ 具有形式
\[
H_t(\omega)=\sum_{k=0}^{n-1}\xi_k(\omega)\mathbf 1_{(t_k,t_{k+1}]}(t),
\]
其中 $0=t_0<t_1<\cdots<t_n=T$，且每个 $\xi_k$ 都是 $\mathscr F_{t_k}$-可测并满足 $\mathbb E[\xi_k^2]<\infty$，则称 $H$ 为\emph{简单适应过程}，并定义其 \emph{Itô 积分} 为
\[
\int_0^T H_t\,dW_t:=\sum_{k=0}^{n-1}\xi_k\bigl(W_{t_{k+1}}-W_{t_k}\bigr).
\]
借助 $L^2(\Omega\times [0,T])$ 完备化，可将该定义延拓到所有平方可积适应过程。
\end{definition}

\begin{proposition}[Itô 等距]\label{prop:ito-isometry}
对任意简单适应过程 $H$，都有
\[
\mathbb E\left[\left(\int_0^T H_t\,dW_t\right)^2\right]
=
\mathbb E\left[\int_0^T |H_t|^2\,dt\right].
\]
因此 Itô 积分是从平方可积适应过程空间到 $L^2(\Omega)$ 的等距映射。
\end{proposition}

\begin{proof}
对定义~\ref{def:ito-integral} 中的简单过程展开平方，得到
\[
\left(\sum_{k=0}^{n-1}\xi_k\Delta W_k\right)^2
=
\sum_{k=0}^{n-1}\xi_k^2(\Delta W_k)^2
+2\sum_{i<j}\xi_i\xi_j\Delta W_i\Delta W_j,
\]
其中 $\Delta W_k:=W_{t_{k+1}}-W_{t_k}$。由于 Brownian motion 具有独立增量，且 $\Delta W_j$ 与 $\mathscr F_{t_j}$ 独立、均值为零，交叉项取期望后消失。于是
\[
\mathbb E\left[\left(\int_0^T H_t\,dW_t\right)^2\right]
=
\sum_{k=0}^{n-1}\mathbb E[\xi_k^2]\,(t_{k+1}-t_k)
=
\mathbb E\left[\int_0^T |H_t|^2\,dt\right].
\]
结论成立。
\end{proof}

\subsection{Itô 公式与最小量级的随机微分方程}

\begin{theorem}[Itô 公式]\label{thm:ito-formula}
设 $W$ 为 Brownian motion，过程 $X$ 满足
\[
X_t=X_0+\int_0^t b_s\,ds+\int_0^t \sigma_s\,dW_s,
\]
其中 $b,\sigma$ 为适应过程，并满足使两项积分有意义的可积条件。若 $f\in C^{1,2}([0,T]\times \mathbb R)$，则对任意 $t\in[0,T]$，
\[
f(t,X_t)=f(0,X_0)
+\int_0^t \partial_t f(s,X_s)\,ds
+\int_0^t \partial_x f(s,X_s)b_s\,ds
+\int_0^t \partial_x f(s,X_s)\sigma_s\,dW_s
+\frac12\int_0^t \partial_{xx}f(s,X_s)\sigma_s^2\,ds.
\]
\end{theorem}

\begin{proof}
先对时间区间做剖分，对每个小区间应用二元 Taylor 展开：
\[
f(t_{k+1},X_{t_{k+1}})-f(t_k,X_{t_k})
\]
分解为时间增量、一阶空间增量和二阶空间增量，再把高阶余项控制到零。关键在于
\[
(\Delta W_k)^2\approx \Delta t_k,\qquad \Delta W_k\Delta t_k,\ (\Delta t_k)^2
\]
在求和极限下都比 $\Delta t_k$ 更高阶，因此只保留二阶项中的二次变差贡献。对随机项与确定项分别取极限，便得到公式。严格论证通常先对简单过程验证，再用 Itô 等距与 $L^2$ 逼近推广。
\end{proof}

\begin{example}[二次函数的 Itô 公式]\label{ex:ito-quadratic}
取 $X_t=W_t$，$f(x)=x^2$。由定理~\ref{thm:ito-formula}，
\[
W_t^2=2\int_0^t W_s\,dW_s+t.
\]
把这个例子放在这里，是为了突出随机微积分与经典微积分最本质的差别：额外出现的 $t$ 项正来自 Brownian motion 的二次变差。
\end{example}

\begin{example}[应用触点：几何 Brownian motion]\label{ex:geometric-bm}
考虑随机微分方程
\[
dS_t=\mu S_t\,dt+\sigma S_t\,dW_t,\qquad S_0>0.
\]
对 $f(x)=\log x$ 应用定理~\ref{thm:ito-formula}，可得
\[
d(\log S_t)=\left(\mu-\frac{\sigma^2}{2}\right)dt+\sigma\,dW_t,
\]
从而
\[
S_t=S_0\exp\left(\left(\mu-\frac{\sigma^2}{2}\right)t+\sigma W_t\right).
\]
把这个例子放在这里，是为了说明第 11 章中的连续时间金融建模和随机控制并不需要额外发明一套新分析，而是直接建立在 Itô 公式之上。
\end{example}

\subsection*{有限维对照}

有限维随机优化通常只显式书写随机向量、协方差矩阵和期望目标；连续时间情形则把“随机向量”替换成路径过程，把“矩阵演化”替换成 Itô 驱动的动力系统。本节的作用，是把这种升维过程中真正新增的结构明确标出：定义~\ref{def:conditional-expectation} 给出信息更新，定义~\ref{def:martingale} 给出公平演化，定义~\ref{def:brownian-motion} 与定义~\ref{def:ito-integral} 给出噪声与积分，而定理~\ref{thm:ito-formula} 则是连续时间随机优化的链式法则。

\subsection*{习题（待补）}

\begin{exercise}[条件期望占位]\label{exr:conditional-expectation-placeholder}
补一题：在离散有限样本空间上显式计算条件期望，并验证命题~\ref{prop:conditional-expectation-basic} 的塔性质。
\end{exercise}

\begin{exercise}[Itô 公式桥接题占位]\label{exr:ito-formula-placeholder}
补一题：从例~\ref{ex:geometric-bm} 出发，验证几何 Brownian motion 的显式解，并说明其如何退化为常微分方程
\[
\dot x_t=\mu x_t
\]
的有限维无噪声版本。
\end{exercise}

